\documentclass[11pt,a4paper]{article}
\usepackage[margin=1in]{geometry}
\usepackage{amsmath,amssymb}
\usepackage{booktabs,tabularx,longtable,array}
\usepackage{enumitem}
\usepackage{titlesec}
\usepackage{fancyhdr}
\usepackage[table]{xcolor}
\usepackage[hidelinks]{hyperref}
\usepackage{tcolorbox}
\tcbuselibrary{skins,breakable}

% Setup page style
\pagestyle{fancy}
\fancyhf{}
\fancyhead[L]{\textbf{Kathmandu University BDS 1st Year}}
\fancyhead[R]{\textbf{Pharmacology Solutions}}
\fancyfoot[C]{\thepage}
\renewcommand{\headrulewidth}{0.6pt}

% Section styling
\titleformat{\section}{\Large\bfseries}{{\thesection}}{1em}{}[\titlerule]
\titleformat{\subsection}{\large\bfseries\color{black!85}}{\thesubsection}{1em}{}
\titleformat{\subsubsection}{\normalsize\bfseries\color{black!75}}{\thesubsubsection}{1em}{}

% Custom boxes
\newtcolorbox{qbox}[2][]{
  colback=gray!5!white,
  colframe=black!45,
  fonttitle=\bfseries,
  title={#2},
  arc=2mm,
  breakable,
  #1
}

\newtcolorbox{ansbox}[1][]{
  colback=white,
  colframe=gray!50!black,
  fonttitle=\bfseries,
  title={Answer},
  arc=1mm,
  breakable,
  #1
}

\begin{document}

\begin{titlepage}
    \centering
    \vspace*{2cm}
    {\Huge \textbf{Kathmandu University}}\\[0.5cm]
    {\LARGE \textbf{Bachelor of Dental Surgery (BDS) 1st Year}}\\[1.5cm]
    {\Huge \textbf{PHARMACOLOGY}}\\[0.8cm]
    {\Large \textbf{Comprehensive Exam Solutions \& Question Bank}}\\[0.5cm]
    {\large Complete Question Papers, Exam-Oriented Answers, Frequency Analysis \& Revision Cheatsheet}\\[2cm]
    \vfill
    {\large \textbf{Compiled from Past Papers: 2015--2022}}\\[0.5cm]
    {\small Kathmandu University School of Medical Sciences (KUSMS) \& Affiliated Colleges}
    \vspace*{1.5cm}
\end{titlepage}

\tableofcontents
\newpage

%========================================
\section{Past Examination Questions and Solutions (2015--2022)}
%========================================

%----------------------------------------
\subsection{Examination: May 20, 2022 (End Semester Exam)}
\textbf{Subject:} Pharmacology \quad \textbf{Marks:} 50 (Section B: 25, Section C: 25)

\subsubsection{Section B [25 Marks]}

\begin{qbox}{Question 1 [2+4=6 Marks]}
What is first pass metabolism? Write in short about factors affecting the absorption of drugs.
\end{qbox}

\begin{ansbox}
\textbf{First-Pass Metabolism (Presystemic Metabolism):}
The phenomenon whereby a drug absorbed from the GIT undergoes significant metabolism in the gut wall and/or liver before reaching the systemic circulation, thereby reducing its bioavailability. Drugs with high first-pass effect (\textgreater{}70\%): morphine, propranolol, nitrates (GTN), lidocaine, verapamil, aspirin. These drugs require higher oral doses or alternative routes (sublingual GTN bypasses hepatic first-pass).

\textbf{Factors Affecting Drug Absorption:}
\begin{enumerate}[leftmargin=*]
    \item \textbf{Physicochemical properties of drug:}
    \begin{itemize}
        \item \textit{Lipid solubility:} Highly lipophilic drugs diffuse readily across cell membranes.
        \item \textit{Ionisation (pKa and pH):} Unionised form is lipid-soluble and absorbable. Henderson-Hasselbalch equation governs ionisation. Weak acids absorbed well from the stomach (low pH); weak bases absorbed better in intestine (high pH).
        \item \textit{Molecular size:} Smaller molecules ($<$500 Da) absorbed better.
        \item \textit{Formulation:} Tablet dissolution rate; coated formulations; modified-release preparations.
    \end{itemize}
    \item \textbf{Route of administration:} IV bypasses absorption; sublingual avoids first-pass.
    \item \textbf{Gastric emptying and intestinal motility:} ↑ motility → faster transit → ↓ absorption time. Food delays gastric emptying → delays absorption of most drugs.
    \item \textbf{Intestinal surface area:} Villi and microvilli; large surface area (200 m²) favours absorption.
    \item \textbf{Splanchnic blood flow:} ↑ blood flow → maintains concentration gradient → ↑ absorption.
    \item \textbf{Drug interactions:} Chelation (e.g., tetracycline + antacids/milk → ↓ absorption); adsorption (charcoal binds drugs).
    \item \textbf{First-pass metabolism:} ↑ hepatic extraction → ↓ bioavailability.
    \item \textbf{P-glycoprotein (MDR-1):} Efflux transporter on intestinal epithelium; pumps drugs back into gut lumen; reduces absorption.
    \item \textbf{Disease states:} Malabsorption, Crohn's disease, gastrectomy → ↓ absorption.
\end{enumerate}
\end{ansbox}

\begin{qbox}{Question 2 [5$\times$2=10 Marks]}
Write briefly on: (a) Prodrug \quad (b) Plasma half life \quad (c) Disadvantages of oral route \quad (d) Essential drug concept \quad (e) Rational drug therapy
\end{qbox}

\begin{ansbox}
\textbf{(a) Prodrug:}
An inactive or poorly active compound that is converted in the body (by enzymatic or chemical reaction) into a pharmacologically active metabolite. Rationale: (i) Improve bioavailability by bypassing first-pass. (ii) Improve stability/solubility/palatability. (iii) Target-specific drug delivery. (iv) Reduce toxicity (NSAID prodrugs). Examples: Enalapril → Enalaprilat (ACE inhibitor, active); Codeine → Morphine (CYP2D6, poor metabolisers don't feel effect); Levodopa → Dopamine (DOPA decarboxylase in brain; dopamine cannot cross BBB); Aspirin → Salicylic acid; Cyclophosphamide → Acrolein.

\textbf{(b) Plasma Half-Life (t$_{1/2}$):}
Time taken for plasma drug concentration to fall by 50\%. Formula: \textbf{t$_{1/2}$ = 0.693 $\times$ Vd / CL} (Vd = volume of distribution; CL = clearance). Clinical importance:
\begin{itemize}[leftmargin=*]
    \item Determines dosing interval (given every 1 t$_{1/2}$ as a rule of thumb).
    \item Steady-state concentration achieved after \textbf{4--5 half-lives}.
    \item Time for drug to be eliminated = 4--5 half-lives.
    \item Short t$_{1/2}$ (e.g., penicillin G: 0.5h) → frequent dosing. Long t$_{1/2}$ (amiodarone: 40--55 days; warfarin: 40h) → accumulation risk.
\end{itemize}

\textbf{(c) Disadvantages of Oral Route:}
\begin{enumerate}[leftmargin=*]
    \item Significant first-pass metabolism reduces bioavailability.
    \item Irregular and incomplete absorption (food, pH, motility affects absorption).
    \item Slow onset of action → unsuitable for emergencies.
    \item Cannot be used in unconscious, uncooperative, or vomiting patients.
    \item GI irritation and ulceration (NSAIDs, iron, potassium chloride).
    \item Some drugs destroyed by gastric acid (insulin, penicillin G benzathine) or digestive enzymes.
\end{enumerate}

\textbf{(d) Essential Drug Concept:}
Essential medicines (WHO, 1977) are those that satisfy the priority healthcare needs of the population. Selected based on: disease prevalence, clinical efficacy, safety, cost-effectiveness, and pharmacokinetic properties. The WHO Essential Medicines List (EML, updated every 2 years) is the global model; countries adapt national formularies. Benefits: improves access, reduces irrational prescribing, lowers healthcare costs.

\textbf{(e) Rational Drug Therapy:}
Prescribing the right drug to the right patient in the right dose by the right route at the right time for the right duration at the lowest cost (WHO/WHO-INRUD definition). Steps: (1) Define diagnosis; (2) Specify therapeutic objective; (3) Select appropriate treatment; (4) Initiate treatment with correct prescription; (5) Inform patient; (6) Monitor treatment and stop/adjust as required. Reduces adverse drug reactions, drug interactions, and healthcare costs.
\end{ansbox}

\begin{qbox}{Question 3 [3$\times$3=9 Marks]}
Write short notes on: (a) Potency and efficacy with illustration \quad (b) Phase I and Phase II metabolic reactions \quad (c) Streptokinase
\end{qbox}

\begin{ansbox}
\textbf{(a) Potency and Efficacy:}

\textbf{Potency:} Concentration (EC$_{50}$) or dose (ED$_{50}$) required to produce 50\% of the maximum response. A drug with \textit{lower} EC$_{50}$ is \textit{more potent} (i.e., needs less drug to produce the same effect).

\textbf{Efficacy (Intrinsic efficacy / E$_{max}$):} The maximum response a drug can produce regardless of dose. A drug with high E$_{max}$ is more efficacious.

Key distinction: A more potent drug is NOT necessarily more efficacious. Example: Morphine has greater efficacy (can relieve severe pain) than aspirin, even though aspirin may be more potent (mg for mg) for mild pain in some contexts. In dose-response curves: two drugs can have the same E$_{max}$ (equal efficacy) but different EC$_{50}$ (different potency); or same EC$_{50}$ but different E$_{max}$ (different efficacy).

Partial agonist: Has lower E$_{max}$ than full agonist even at maximal receptor occupancy; can act as functional antagonist in presence of full agonist.

\textbf{(b) Phase I and Phase II Metabolic Reactions:}

\textbf{Phase I (``Functionalisation''):}
Introduce or reveal a functional group (-OH, -NH$_2$, -COOH, -SH). Make drug more polar. Mediated mainly by \textbf{cytochrome P450 (CYP)} enzymes in the liver (and gut, lung). Most important isoforms: CYP3A4 (60\% of drugs), CYP2D6, CYP2C9, CYP2C19. Types: Oxidation (hydroxylation, N-/O-dealkylation, S-oxidation), Reduction (nitro-reduction), Hydrolysis (ester, amide hydrolysis). Phase I product may still be pharmacologically active.

\textbf{Phase II (``Conjugation / Biosynthetic''):}
Attach an endogenous polar molecule to Phase I product (or occasionally to parent drug directly) → makes drug water-soluble and pharmacologically inactive → urinary or biliary excretion.
\begin{itemize}[leftmargin=*]
    \item \textbf{Glucuronidation} (most common; UDP-glucuronyltransferases; liver, kidney, gut).
    \item \textbf{Sulfation} (phenols, steroids; saturable).
    \item \textbf{Acetylation} (N-acetyltransferases; ``fast'' vs. ``slow'' acetylators; isoniazid — slow acetylators more prone to neuropathy).
    \item \textbf{Methylation} (COMT for catecholamines; TPMT for 6-mercaptopurine).
    \item \textbf{Glutathione conjugation} (hepatoprotective; N-acetylcysteine replenishes glutathione in paracetamol overdose).
    \item \textbf{Glycine conjugation} (bile acids, salicylates).
\end{itemize}

\textbf{(c) Streptokinase:}
Thrombolytic (fibrinolytic) agent derived from Group C $\beta$-haemolytic streptococci. Mechanism of action: Streptokinase binds to plasminogen → forms streptokinase-plasminogen complex → acts as a plasminogen activator → converts free plasminogen → \textbf{plasmin} → plasmin degrades fibrin clot (and fibrinogen, factors V, VIII). Not fibrin-selective (``lytic state'' = systemic fibrinogenolysis). Uses: Acute STEMI (within 12h, no PCI facility), acute massive PE, acute DVT, peripheral arterial thrombosis. Adverse effects: Haemorrhage (major; ICH), \textbf{anaphylaxis/allergic reactions} (streptococcal antigen, antibodies in re-use), fever, hypotension. Contraindicated in: recent surgery (\textless{}2 weeks), prior stroke, active GI bleeding, severe uncontrolled hypertension.
\end{ansbox}

\subsubsection{Section C [25 Marks]}

\begin{qbox}{Question 4 [3+2+1=6 Marks]}
Classify anticoagulant drugs. Write down the mechanism of action of warfarin. Mention its important therapeutic uses.
\end{qbox}

\begin{ansbox}
\textbf{Classification of Anticoagulant Drugs:}

\textbf{A. Parenteral Anticoagulants:}
\begin{enumerate}[leftmargin=*]
    \item \textbf{Unfractionated Heparin (UFH):} Indirect thrombin inhibitor; potentiates antithrombin III (AT-III) → inhibits thrombin (IIa) and Xa (equally). IV/SC. Monitoring: APTT. Antidote: Protamine sulfate.
    \item \textbf{Low Molecular Weight Heparins (LMWH):} Enoxaparin, Dalteparin. Mainly anti-Xa activity; more predictable pharmacokinetics; SC; no APTT monitoring required.
    \item \textbf{Fondaparinux:} Synthetic pentasaccharide; selective factor Xa inhibitor; no anti-IIa activity. SC.
    \item \textbf{Direct Thrombin Inhibitors (parenteral):} Bivalirudin, Argatroban, Lepirudin.
\end{enumerate}

\textbf{B. Oral Anticoagulants:}
\begin{enumerate}[leftmargin=*]
    \item \textbf{Vitamin K antagonists:} Warfarin, Acenocoumarol, Phenindione.
    \item \textbf{Direct Oral Anticoagulants (DOACs):}
    \begin{itemize}
        \item Direct Xa inhibitors: Rivaroxaban, Apixaban, Edoxaban.
        \item Direct IIa (thrombin) inhibitor: Dabigatran.
    \end{itemize}
\end{enumerate}

\textbf{Mechanism of Action of Warfarin:}
\begin{itemize}[leftmargin=*]
    \item Warfarin is a coumarin derivative; orally active.
    \item Inhibits \textbf{Vitamin K Epoxide Reductase (VKOR)} → prevents regeneration of active reduced Vitamin K (Vit K$_1$ hydroquinone).
    \item Active Vit K is essential cofactor for $\gamma$-glutamyl carboxylase enzyme → adds $\gamma$-carboxyl groups to glutamate residues of \textbf{coagulation factors II, VII, IX, X} (``1972'' mnemonic) and anticoagulant proteins \textbf{C and S}.
    \item Without carboxylation, these factors cannot bind $\text{Ca}^{2+}$ via phospholipid membrane (PIVKA = Proteins Induced by Vitamin K Absence) → cannot participate in coagulation cascade.
    \item \textbf{Delayed onset:} 8--12 hours (existing factors must clear); full effect after 3--5 days. Factor VII has shortest half-life → PT prolongs first.
    \item Monitoring: \textbf{INR (International Normalised Ratio)} = target 2--3 for most indications (2.5--3.5 for mechanical heart valves).
\end{itemize}

\textbf{Important Therapeutic Uses of Warfarin:}
\begin{itemize}[leftmargin=*]
    \item Deep Vein Thrombosis (DVT) and Pulmonary Embolism treatment and prophylaxis.
    \item Atrial Fibrillation (AF) -- prevention of cardioembolic stroke.
    \item Mechanical prosthetic heart valves (anticoagulation of choice).
    \item Prevention of recurrent MI and stroke.
    \item Antiphospholipid antibody syndrome.
\end{itemize}
\end{ansbox}

\begin{qbox}{Question 5 [3+2=5 Marks]}
Write the classification of adrenergic drugs on the basis of their therapeutic effects with at least one example of each. Write down the therapeutic uses of Prazosin.
\end{qbox}

\begin{ansbox}
\textbf{Classification of Adrenergic Drugs (by Therapeutic Effect):}

\textbf{A. Adrenergic Agonists (Sympathomimetics):}
\begin{enumerate}[leftmargin=*]
    \item \textbf{Vasopressors/Cardiotonic:}
    \begin{itemize}
        \item $\alpha_1 + \beta_1$ agonist: Noradrenaline (Norepinephrine) -- vasopressor in septic shock.
        \item $\beta_1$ agonist: Dobutamine -- acute heart failure.
    \end{itemize}
    \item \textbf{Bronchodilators ($\beta_2$ agonists):}
    \begin{itemize}
        \item Short-acting (SABA): Salbutamol (albuterol) -- acute bronchospasm.
        \item Long-acting (LABA): Salmeterol, Formoterol -- COPD/asthma maintenance.
    \end{itemize}
    \item \textbf{Nasal Decongestants ($\alpha_1$ agonists):} Xylometazoline, Oxymetazoline.
    \item \textbf{Anaphylaxis Management:} Adrenaline (Epinephrine) -- $\alpha_1+\beta_1+\beta_2$ -- drug of choice in anaphylaxis; also causes bronchodilation and vasoconstriction.
    \item \textbf{Uterine relaxants (tocolytics, $\beta_2$):} Ritodrine, Terbutaline.
    \item \textbf{CNS stimulants:} Amphetamines (indirect-acting; release NA and DA).
\end{enumerate}

\textbf{B. Adrenergic Antagonists (Sympatholytics):}
\begin{enumerate}[leftmargin=*]
    \item \textbf{$\alpha_1$ blockers (antihypertensives/BPH):} Prazosin, Tamsulosin ($\alpha_{1A}$ selective for BPH), Doxazosin, Terazosin.
    \item \textbf{$\alpha_2$ blockers:} Yohimbine (experimental).
    \item \textbf{Non-selective $\alpha$ blockers (phaeochromocytoma):} Phenoxybenzamine (irreversible), Phentolamine.
    \item \textbf{$\beta$ blockers:}
    \begin{itemize}
        \item Non-selective: Propranolol ($\beta_1 + \beta_2$) -- hypertension, angina, arrhythmia.
        \item Selective ($\beta_1$): Atenolol, Metoprolol, Bisoprolol -- hypertension, HF, post-MI.
        \item With $\alpha$-blocking: Labetalol, Carvedilol -- hypertensive crisis, HF.
    \end{itemize}
\end{enumerate}

\textbf{Therapeutic Uses of Prazosin ($\alpha_1$-adrenergic blocker):}
\begin{enumerate}[leftmargin=*]
    \item \textbf{Hypertension:} First-dose hypotension effect (``prazosin syncope'') -- give at bedtime; overcomes by starting with low dose. Dilates arterioles and veins.
    \item \textbf{Benign Prostatic Hyperplasia (BPH):} $\alpha_{1A}$ receptors on prostatic smooth muscle → relaxation → improved urinary flow.
    \item \textbf{Raynaud's phenomenon:} Vasodilation in peripheral arteries.
    \item \textbf{Congestive Heart Failure (adjunct):} Reduces preload and afterload.
    \item \textbf{Phaeochromocytoma (preoperative):} Controls hypertensive crises before surgery.
    \item \textbf{PTSD-related nightmares:} Prazosin shown to reduce nightmares in PTSD.
\end{enumerate}
\end{ansbox}

\begin{qbox}{Question 6 [3+1+1=5 Marks]}
Explain the mechanism of action of Chloroquine. Write down its important therapeutic uses and adverse effects.
\end{qbox}

\begin{ansbox}
\textbf{Mechanism of Action of Chloroquine:}

Chloroquine is a 4-aminoquinoline antimalarial.

\begin{enumerate}[leftmargin=*]
    \item \textbf{Primary mechanism (blood schizontocide):}
    \begin{itemize}
        \item Parasitised RBCs digest haemoglobin in the food vacuole → release haem (ferriprotoporphyrin IX = FPIX).
        \item Free FPIX is toxic to the parasite (inhibits enzymes, damages membranes).
        \item Normally the parasite \textbf{polymerises FPIX into inert haemozoin (malaria pigment)}.
        \item Chloroquine (weak base) concentrates in the acidic food vacuole (pH 5.5) by ``ion trapping''; \textbf{inhibits haem polymerisation} → free haem accumulates → toxic to parasite.
    \end{itemize}
    \item \textbf{Additional mechanism (for SLE/RA):} Chloroquine inhibits lysosomal enzymes, reduces antigen processing (AP via MHC class II), reduces cytokine production (IL-1, IL-6, TNF-$\alpha$), and raises lysosomal pH → immunomodulatory.
\end{enumerate}

\textbf{Chloroquine is active only against:} Asexual blood stages (ring, trophozoite, schizont) of \textit{P. vivax}, \textit{P. malariae}, \textit{P. ovale} and sensitive \textit{P. falciparum} (widespread resistance now). Does not affect liver stages.

\textbf{Important Therapeutic Uses:}
\begin{itemize}[leftmargin=*]
    \item \textbf{Malaria treatment:} Drug of choice for sensitive \textit{P. vivax}, \textit{P. malariae}, \textit{P. ovale}, and chloroquine-sensitive \textit{P. falciparum}.
    \item \textbf{Malaria prophylaxis:} In areas without chloroquine resistance.
    \item \textbf{Rheumatoid arthritis (RA) -- DMARD:} Hydroxychloroquine more commonly used.
    \item \textbf{Systemic Lupus Erythematosus (SLE):} Reduces disease flares.
    \item \textbf{Amoebic hepatic abscess} (extra-intestinal amoebiasis): Chloroquine concentrates in liver.
\end{itemize}

\textbf{Adverse Effects:}
\begin{itemize}[leftmargin=*]
    \item \textbf{GI:} Nausea, vomiting, abdominal cramps, diarrhoea (most common).
    \item \textbf{Ocular (most important -- dose-related):} Retinopathy (deposits in retinal pigment epithelium → ``Bull's eye'' maculopathy, visual field defects, colour vision disturbances). Corneal deposits (reversible). Annual ophthalmologic screening required for prolonged use.
    \item \textbf{Dermatological:} Pruritus (especially in dark-skinned patients), photosensitivity, pigmentation changes.
    \item \textbf{CNS:} Headache, dizziness, psychosis (rare).
    \item \textbf{Cardiac:} QT prolongation, cardiomyopathy (rare but serious with high doses).
    \item \textbf{Haematological:} Haemolysis in G6PD-deficient patients (less common than primaquine).
    \item \textbf{Acute overdose:} Life-threatening: hypotension, cardiac arrhythmias, respiratory arrest; fatal if $>$5g in adults.
\end{itemize}
\end{ansbox}

\begin{qbox}{Question 7 [3$\times$3=9 Marks]}
Write short notes on: (a) Organophosphorous poisoning and its management \quad (b) Therapeutic uses of Atropine \quad (c) Therapeutic uses of Propranolol
\end{qbox}

\begin{ansbox}
\textbf{(a) Organophosphorus (OP) Poisoning and Management:}

OP compounds (e.g., Malathion, Parathion, Diazinon; also nerve agents Sarin, Tabun) are irreversible inhibitors of acetylcholinesterase (AChE).

\textbf{Mechanism:} OP binds covalently and irreversibly to the serine hydroxyl at the active site of AChE → \textbf{acetylcholine (ACh) accumulates} → overstimulation of both muscarinic and nicotinic receptors.

\textbf{Clinical Features (SLUD mnemonic for muscarinic):}
\begin{itemize}[leftmargin=*]
    \item \textbf{Muscarinic (DUMBELS):} Diarrhoea/Defecation, Urination, Miosis (pin-point pupils), Bradycardia, Excessive secretions (lacrimation, salivation, bronchorrhoea, sweating), Lacrimation, Vomiting.
    \item \textbf{Nicotinic:} Muscle fasciculations, weakness, paralysis (respiratory muscles → respiratory failure).
    \item \textbf{CNS:} Anxiety, convulsions, coma.
\end{itemize}

\textbf{Management:}
\begin{enumerate}[leftmargin=*]
    \item \textbf{Decontamination:} Remove from exposure; remove clothing; wash skin with soap and water.
    \item \textbf{Supportive care:} Airway management (intubation), oxygen, IV fluids, control seizures (diazepam).
    \item \textbf{Specific antidotes:}
    \begin{itemize}
        \item \textbf{Atropine:} Competitive muscarinic antagonist; given IV in large, repeated doses (2--4 mg every 10--15 min) until secretions dry up (atropinisation). Reverses muscarinic effects only; not nicotinic.
        \item \textbf{Pralidoxime (2-PAM, oximes):} Reactivates AChE if given \textbf{before ``ageing''} (irreversible covalent bonding becomes complete; within 24--48h for most OP). Reverses both muscarinic AND nicotinic effects; particularly helpful for muscle paralysis.
    \end{itemize}
\end{enumerate}

\textbf{(b) Therapeutic Uses of Atropine:}

Atropine is a competitive non-selective muscarinic receptor antagonist (anticholinergic).
\begin{enumerate}[leftmargin=*]
    \item \textbf{Preanaesthetic medication:} Reduces airway and salivary secretions (antisialogogue), prevents vagal bradycardia during anaesthesia.
    \item \textbf{Bradyarrhythmias:} Drug of choice for acute sinus bradycardia with haemodynamic compromise, symptomatic bradycardia, and AV block (particularly during inferior MI).
    \item \textbf{Antidote for organophosphorus (OP) poisoning} and muscarinic excess (as above).
    \item \textbf{Ophthalmology:} Mydriasis and cycloplegia for fundal examination; refraction in children; uveitis (prevents synechiae).
    \item \textbf{Antispasmodic:} Peptic ulcer (↓ gastric acid secretion), irritable bowel syndrome, renal/biliary colic (smooth muscle relaxation).
    \item \textbf{Bronchospasm:} Ipratropium bromide (quaternary derivative, inhaled) for COPD and asthma.
    \item \textbf{Motion sickness and vertigo} (scopolamine/hyoscine transdermal patch).
    \item \textbf{Antidote for poisoning by muscarinic agonists:} Pilocarpine, muscarine, physostigmine.
\end{enumerate}

\textbf{(c) Therapeutic Uses of Propranolol:}

Propranolol is a non-selective $\beta$-adrenergic receptor antagonist ($\beta_1 + \beta_2$ blockade).
\begin{enumerate}[leftmargin=*]
    \item \textbf{Hypertension:} First-line in young patients; reduces CO and renin release.
    \item \textbf{Angina pectoris:} Reduces myocardial oxygen demand (↓ HR, ↓ contractility, ↓ BP).
    \item \textbf{Cardiac arrhythmias:} Supraventricular tachycardia, atrial fibrillation rate control; Wolff-Parkinson-White syndrome.
    \item \textbf{Post-Myocardial Infarction:} Reduces mortality (reduces infarct size, prevents reinfarction, antiarrhythmic).
    \item \textbf{Heart Failure (with caution):} Carvedilol/metoprolol preferred; propranolol reduces mortality in stable HF.
    \item \textbf{Thyrotoxicosis / Thyroid storm:} Blocks adrenergic manifestations (tachycardia, tremor, sweating, anxiety); also inhibits peripheral conversion of T4 → T3.
    \item \textbf{Phaeochromocytoma (after $\alpha$-blockade):} Controls tachycardia; \textbf{must not give $\beta$-blocker before $\alpha$-blocker} (causes paradoxical hypertension).
    \item \textbf{Migraine prophylaxis:} Reduces frequency and severity of attacks.
    \item \textbf{Essential tremor:} $\beta_2$ receptor blockade in peripheral muscles.
    \item \textbf{Portal hypertension:} Reduces splanchnic blood flow; prophylaxis of variceal bleeding.
    \item \textbf{Anxiety (situational):} Peripheral $\beta$ blockade reduces somatic symptoms (palpitations, tremor).
    \item \textbf{Hypertrophic obstructive cardiomyopathy (HOCM):} Reduces outflow obstruction.
\end{enumerate}

\textbf{Important Contraindications of Propranolol:} Bronchial asthma / COPD ($\beta_2$ blockade causes bronchospasm); acute decompensated heart failure; AV heart block ($>$1st degree); Prinzmetal's angina.
\end{ansbox}

%----------------------------------------
\subsection{Examination: March 15, 2022}

\subsubsection{Section B [25 Marks]}

\begin{qbox}{Question [2+1+2+1=6 Marks]}
Classify anticholinergic drugs. Write briefly mechanism of action, indications and adverse effects of the prototype anticholinergic drug.
\end{qbox}

\begin{ansbox}
\textbf{Classification of Anticholinergic (Antimuscarinic) Drugs:}
\begin{enumerate}[leftmargin=*]
    \item \textbf{Tertiary amines (lipid-soluble, cross BBB):}
    \begin{itemize}
        \item Natural alkaloids: Atropine, Scopolamine (Hyoscine).
        \item Synthetic: Benzhexol (Trihexyphenidyl, Parkinson's); Dicyclomine (IBS); Oxybutynin (OAB); Pirenzepine (selective M$_1$, peptic ulcer).
    \end{itemize}
    \item \textbf{Quaternary ammonium compounds (poorly absorbed, limited CNS entry):}
    \begin{itemize}
        \item Ipratropium bromide (COPD/asthma inhalation), Tiotropium (COPD), Glycopyrrolate (anaesthesia).
    \end{itemize}
\end{enumerate}

\textbf{Prototype: Atropine -- Mechanism, Indications, Adverse Effects:}

\textbf{Mechanism:} Competitive antagonist at muscarinic receptors (M$_1$--M$_5$); non-selective; blocks all actions of ACh/cholinergic drugs at muscarinic sites. No effect on nicotinic receptors. Effect is dose-dependent.

\textbf{Indications:} (Listed in Answer 7b above -- preanaesthetic, bradycardia, OP poisoning, ophthalmology, antispasmodic, bronchospasm, motion sickness.)

\textbf{Adverse Effects:} (Anticholinergic syndrome, mnemonic: ``Dry as a bone, Blind as a bat, Red as a beet, Hot as a hare, Mad as a hatter, Full as a flask'')
\begin{itemize}[leftmargin=*]
    \item \textbf{Dry mouth} (xerostomia -- reduced salivation); Blurred vision (cycloplegia, mydriasis); Photophobia.
    \item \textbf{Urinary retention} (especially in elderly men with BPH).
    \item \textbf{Constipation} (↓ gut motility).
    \item \textbf{Tachycardia} (vagal blockade at heart).
    \item \textbf{Flushing and dry skin} (sweat gland inhibition → hyperthermia).
    \item \textbf{CNS effects:} Confusion, restlessness, hallucinations (esp. scopolamine and central-acting antimuscarinics).
    \item \textbf{Acute angle-closure glaucoma} (pupillary dilation blocks angle drainage → raised IOP). Contraindicated in narrow-angle glaucoma.
\end{itemize}
\end{ansbox}

%========================================
\subsection{Examination: October 31, 2021}

\begin{qbox}{Question [4+2=6 Marks]}
List different routes of drug administration with an example of each. Write two advantages and two disadvantages of sublingual route.
\end{qbox}

\begin{ansbox}
\textbf{Routes of Drug Administration:}
\begin{center}
\begin{tabular}{lll}
\toprule
\textbf{Route} & \textbf{Sub-types} & \textbf{Example} \\
\midrule
Enteral & Oral & Paracetamol tablet \\
 & Sublingual (SL) & GTN (glyceryl trinitrate) \\
 & Buccal & Fentanyl buccal film \\
 & Rectal & Diazepam suppository \\
Parenteral & Intravenous (IV) & Morphine injection \\
 & Intramuscular (IM) & Vaccines, diclofenac \\
 & Subcutaneous (SC) & Insulin, LMWH \\
 & Intradermal (ID) & Mantoux test \\
 & Intrathecal & Spinal anaesthesia \\
Inhalation & -- & Salbutamol MDI \\
Topical/Transdermal & Skin, eye, ear, nose & Clotrimazole cream; GTN patch \\
\bottomrule
\end{tabular}
\end{center}

\textbf{Advantages of Sublingual (SL) Route:}
\begin{enumerate}[leftmargin=*]
    \item \textbf{Rapid onset:} Drug is absorbed directly from sublingual mucosa into systemic circulation via sublingual veins → superior vena cava → heart; bypass of hepatic first-pass metabolism → faster onset than oral route. GTN relieves angina within 1--2 minutes.
    \item \textbf{Avoids first-pass effect:} Higher bioavailability for drugs with high hepatic extraction ratio (e.g., GTN, buprenorphine, testosterone, fentanyl).
\end{enumerate}

\textbf{Disadvantages of Sublingual Route:}
\begin{enumerate}[leftmargin=*]
    \item \textbf{Limited drug absorption:} Only suitable for lipid-soluble, non-irritating drugs that can be absorbed through the thin mucosa. Drugs with high molecular weight, poor lipid solubility, or ionised form cannot be administered SL.
    \item \textbf{Patient cooperation required:} Patient must hold tablet under the tongue without swallowing or speaking; difficult in uncooperative, unconscious, or vomiting patients.
\end{enumerate}
\end{ansbox}

%========================================
\section{Frequency Analysis of Examination Topics}
%========================================

\subsection{Most Frequently Repeated Topics (80\%+ of Papers)}
\begin{itemize}[leftmargin=*]
    \item \textbf{Drug Absorption / Bioavailability / Pharmacokinetics (90\%):} Factors affecting absorption (lipid solubility, pKa, first-pass, P-gp), bioavailability, first-pass metabolism, routes of administration.
    \item \textbf{Beta Blockers (85\%):} Classification (selective vs. non-selective); propranolol mechanism and therapeutic uses; atenolol advantages; carvedilol.
    \item \textbf{Anticoagulants --- classification and warfarin mechanism (85\%):} Heparin (UFH, LMWH, mechanism, antidote: protamine), warfarin (VKOR inhibitor, INR monitoring, drug interactions), DOACs.
\end{itemize}

\subsection{Frequently Repeated Topics (60\%--79\% of Papers)}
\begin{itemize}[leftmargin=*]
    \item \textbf{Digoxin --- mechanism of action (75\%):} Na/K ATPase inhibition → ↑ intracellular Na → ↑ intracellular Ca → positive inotropy; also vagomimetic (↓ HR, ↓ AV conduction). Uses: HF, AF.
    \item \textbf{Anticholinergic drugs / Atropine (70\%):} Classification, mechanism, therapeutic uses, adverse effects, contraindications.
    \item \textbf{Antimalarial drugs --- Chloroquine (65\%):} Mechanism (haem polymerisation inhibition), uses, adverse effects (retinopathy); primaquine (radical cure for vivax), mefloquine.
    \item \textbf{Antitubercular drugs (65\%):} First-line HRZE, mechanisms, adverse effects, DOT, short-course chemotherapy.
    \item \textbf{Antifungal drugs (60\%):} Amphotericin B (binds ergosterol); azoles (lanosterol → ergosterol synthesis inhibition); clotrimazole (oral thrush); fluconazole.
    \item \textbf{Routes of Administration (60\%):} Routes, examples, advantages/disadvantages of oral, IV, sublingual, IM, SC, inhalation.
\end{itemize}

\subsection{Moderately Repeated Topics (40\%--59\% of Papers)}
\begin{itemize}[leftmargin=*]
    \item \textbf{Antianginal drugs (55\%):} Nitrates (cGMP → MLCK phosphorylation → smooth muscle relaxation → venodilation → ↓ preload; also coronary dilation); $\beta$-blockers; CCBs.
    \item \textbf{Cephalosporins (50\%):} Generations 1--4, mechanism (inhibit transpeptidase/PBP), adverse effects (cross-sensitivity with penicillin, nephrotoxicity).
    \item \textbf{Fluoroquinolones (50\%):} Mechanism (DNA gyrase/topoisomerase IV inhibition); ciprofloxacin uses; adverse effects (tendinopathy, QT prolongation, avoid in children/pregnancy).
    \item \textbf{Penicillins (50\%):} Classification, mechanism, amoxicillin + clavulanic acid rationale.
    \item \textbf{Bronchodilators / Asthma drugs (45\%):} Salbutamol ($\beta_2$ agonist), ipratropium, aminophylline (methylxanthine; PDE inhibitor → cAMP $\uparrow$; narrow TI).
    \item \textbf{Organophosphorus poisoning (40\%):} Mechanism, features, atropine + pralidoxime management.
    \item \textbf{Heparin mechanism and uses (40\%):} AT-III activation, monitoring (APTT), protamine antidote, LMWHs.
\end{itemize}

\subsection{Rarely Repeated Topics (20\%--39\% of Papers)}
\begin{itemize}[leftmargin=*]
    \item \textbf{Streptokinase (35\%):} Mechanism, uses (STEMI, PE), adverse effects (anaphylaxis, bleeding).
    \item \textbf{Calcium Channel Blockers (30\%):} Amlodipine (dihydropyridine, vascular smooth muscle); Verapamil/Diltiazem (cardiac; AV block risk).
    \item \textbf{Drug metabolism Phases I and II (30\%):} CYP450, glucuronidation, acetylation (fast/slow).
    \item \textbf{Adrenergic drugs classification (25\%):} Prazosin therapeutic uses.
    \item \textbf{Anticancer / Antiretroviral drugs (20\%):} Principles of chemotherapy; NRTIs, NNRTIs, PIs in HIV.
\end{itemize}

\subsection{Unique / One-Time Questions ($<20\%$ of Papers)}
\begin{itemize}[leftmargin=*]
    \item Hypolipidaemic drugs (statins -- HMG-CoA reductase inhibitors).
    \item Aminoglycosides mechanism, uses, adverse effects (nephrotoxicity, ototoxicity).
    \item Antihistamines (H$_1$-receptor antagonists) -- uses, generations, adverse effects.
    \item Acyclovir mechanism (DNA polymerase inhibition) for herpes infections.
    \item Erythropoietin in anaemia management.
    \item Drug therapy of megaloblastic anaemia (cyanocobalamin, folic acid).
\end{itemize}

\newpage
\appendix
\section{Pharmacology Quick Revision Cheatsheet}

\subsection{Autonomic Pharmacology Summary}
\begin{itemize}[leftmargin=*]
    \item \textbf{Cholinergic drugs:} Direct agonists: Pilocarpine (M), Bethanechol (M, urinary retention), Carbachol; Indirect (anticholinesterases): Neostigmine (quaternary, MG, reversal of NMB), Physostigmine (tertiary, glaucoma, antidote), Organophosphates (irreversible).
    \item \textbf{Anticholinergic syndrome (``dry as a bone''):} Tachycardia, dry mouth/skin, urinary retention, constipation, mydriasis, hyperthermia, confusion. Antidote: Physostigmine (crosses BBB).
    \item \textbf{Sympathomimetics:} Adrenaline = $\alpha_1 + \beta_1 + \beta_2$ (anaphylaxis, cardiac arrest); Noradrenaline = $\alpha_1 + \beta_1$ (vasopressor); Dobutamine = $\beta_1$ (acute HF); Salbutamol = $\beta_2$ (asthma); Phenylephrine = $\alpha_1$ (nasal decongestant, pupil dilation).
    \item \textbf{Propranolol:} Non-selective $\beta$-blocker; CI in asthma; uses: HTN, angina, MI, arrhythmia, thyroid storm, migraine prophylaxis, essential tremor, portal hypertension.
    \item \textbf{Prazosin:} $\alpha_1$-blocker; first-dose hypotension; uses: HTN, BPH.
\end{itemize}

\subsection{Cardiovascular Pharmacology Summary}
\begin{itemize}[leftmargin=*]
    \item \textbf{Digoxin MOA:} Inhibits Na$^+$/K$^+$-ATPase $\rightarrow$ ↑ intracellular Na$^+$ $\rightarrow$ ↑ intracellular Ca$^{2+}$ (via Na-Ca exchanger reversal) $\rightarrow$ ↑ contractility. Also vagotonic (AV node: ↓ conduction/HR -- useful in AF). Narrow therapeutic index; toxicity: nausea, xanthopsia (yellow vision), arrhythmias. Antidote: digoxin-specific Fab fragments.
    \item \textbf{Nitrates MOA:} Organic nitrates metabolised to NO → activates guanylate cyclase → ↑ cGMP → activates PKG → dephosphorylates myosin light chain → smooth muscle relaxation → venodilation (mainly) → ↓ preload → ↓ cardiac work. Also dilates coronary arteries. Tolerance: take nitrate-free interval.
    \item \textbf{Warfarin:} Inhibits VKOR → blocks carboxylation of factors II, VII, IX, X; protein C, S. Monitoring: INR. Antidote: Vitamin K$_1$ (phytomenadione) for minor bleeding; Fresh Frozen Plasma (FFP) + Vit K for major bleeding. CYP2C9 metabolism → many drug interactions.
    \item \textbf{Heparin vs. Warfarin:} Heparin -- immediate, parenteral, APTT monitoring, protamine antidote, safe in pregnancy (doesn't cross placenta). Warfarin -- delayed onset, oral, INR monitoring, Vit K antidote, teratogenic (crosses placenta).
\end{itemize}

\subsection{Antimicrobial Pharmacology Summary}
\begin{itemize}[leftmargin=*]
    \item \textbf{$\beta$-lactam ring:} Nucleus of penicillins, cephalosporins, carbapenems, monobactams. Inhibits transpeptidases (PBPs) → prevents peptidoglycan cross-linking → osmotic lysis.
    \item \textbf{Penicillin allergy + Clavulanic acid rationale:} $\beta$-lactamase-producing bacteria (S. aureus, H. influenzae, E. coli) destroy amoxicillin. Clavulanic acid is a $\beta$-lactamase inhibitor (suicide inhibitor) → protects amoxicillin.
    \item \textbf{Aminoglycosides:} Gentamicin, Tobramycin, Amikacin. Inhibit 30S ribosome → misreading of mRNA. Synergistic with $\beta$-lactams. Toxicity: \textbf{nephrotoxicity} (dose-related, proximal tubule) and \textbf{ototoxicity} (vestibular and cochlear; irreversible). Monitor peak and trough levels.
    \item \textbf{Antifungal MOA:} Amphotericin B: binds ergosterol → pore formation → K$^+$ efflux → cell death (fungicidal). Azoles (fluconazole, itraconazole, voriconazole): inhibit CYP51 (lanosterol → ergosterol) → disrupts membrane. Echinocandins (caspofungin): inhibit 1,3-$\beta$-D-glucan synthase → disrupts cell wall (fungicidal for Candida).
    \item \textbf{Anti-TB (First-line HRZE):} H = Isoniazid (inhibits mycolic acid synthesis); R = Rifampicin (inhibits RNA polymerase $\beta$-subunit); Z = Pyrazinamide (disrupts membrane potential, active in acid pH); E = Ethambutol (inhibits arabinosyl transferase → arabinogalactan synthesis). Common side effects: INH -- peripheral neuropathy (prevent with pyridoxine Vit B6), hepatotoxicity; Rifampicin -- orange urine/body fluids, enzyme inducer; Pyrazinamide -- hyperuricaemia, hepatotoxicity; Ethambutol -- optic neuritis (check colour vision monthly).
\end{itemize}

\vspace{1cm}
\begin{center}
\textbf{--- End of Pharmacology Solutions ---}
\end{center}

\end{document}
